\chapter{Weak Solutions, Part I}

Let $\Omega\subset\R^n$ be a domain. A function $u\in H^1(\Omega)$ is a weak
solution of
\[
-D_j(a_{ij}(x)D_i u)+c(x)u=f(x)
\]
if
\[
\int_\Omega \bigl(a_{ij}D_i uD_j\varphi+cu\varphi\bigr)
=\int_\Omega f\varphi
\qquad\text{for every }\varphi\in H^1_0(\Omega).
\]
Unless stated otherwise, $a_{ij}\in L^\infty(\Omega)$ and, for some
$\lambda>0$,
\[
a_{ij}(x)\xi_i\xi_j\geq\lambda|\xi|^2
\qquad(x\in\Omega,\ \xi\in\R^n),
\]
while $c\in L^{n/2}(\Omega)$ and
$f\in L^{2n/(n+2)}(\Omega)$.

\section{Growth of Local Integrals}

For $B_r(x)\subset\Omega$ and $u\in L^1(\Omega)$, write
\[
u_{x,r}=\frac{1}{|B_r(x)|}\int_{B_r(x)}u.
\]

\begin{theorem}\label{hanlin:weaksol-part1-campanato}\dcref{Theorem 3.1}\group{ch03-weak-solutions-part-i}
\source{han-lin-lecture-notes:page-0058,han-lin-lecture-notes:page-0059}
Suppose $u\in L^2(\Omega)$ and, for some $\alpha\in(0,1)$,
\[
\int_{B_r(x)}|u-u_{x,r}|^2\leq M^2r^{n+2\alpha}
\qquad\text{whenever }B_r(x)\subset\Omega.
\]
Then $u\in C^\alpha(\Omega)$, and for every $\Omega'\Subset\Omega$,
\[
\sup_{\Omega'}|u|+
\sup_{\substack{x,y\in\Omega'\\x\ne y}}
\frac{|u(x)-u(y)|}{|x-y|^\alpha}
\leq c\bigl(M+\|u\|_{L^2(\Omega)}\bigr),
\]
where $c$ depends only on $n,\alpha,\Omega$, and $\Omega'$.
\end{theorem}
\begin{proof}
\sketch Put $R_0=\dist(\Omega',\partial\Omega)$. Comparing averages at
$0<r_1<r_2\leq R_0$ gives
\[
|u_{x,r_1}-u_{x,r_2}|^2
\leq c(n)M^2r_1^{-n}
\bigl(r_1^{n+2\alpha}+r_2^{n+2\alpha}\bigr).
\]
Summation over dyadic radii shows that the averages converge to a representative
$\widehat u$ and that
\[
|u_{x,r}-\widehat u(x)|\leq c(n,\alpha)Mr^\alpha.
\]
Lebesgue differentiation identifies $\widehat u=u$ almost everywhere; uniform
convergence of the averages on $\Omega'$ makes this representative continuous.
Taking $r=R_0$ yields
\[
\sup_{\Omega'}|u|
\leq c\bigl(MR_0^\alpha+\|u\|_{L^2(\Omega)}\bigr).
\]
If $R=|x-y|<R_0/2$, average over
$B_{2R}(x)\cap B_{2R}(y)$ to obtain
\[
|u_{x,2R}-u_{y,2R}|^2\leq c(n,\alpha)M^2R^{2\alpha},
\]
and hence $|u(x)-u(y)|\leq cM|x-y|^\alpha$. The preceding supremum
bound handles $|x-y|\geq R_0/2$.
\end{proof}

\begin{corollary}\label{hanlin:weaksol-part1-morrey}\dcref{Corollary 3.2}\group{ch03-weak-solutions-part-i}
\source{han-lin-lecture-notes:page-0059}
Suppose $u\in H^1_{\mathrm{loc}}(\Omega)$ and, for some
$\alpha\in(0,1)$,
\[
\int_{B_r(x)}|Du|^2\leq M^2r^{n-2+2\alpha}
\qquad\text{whenever }B_r(x)\subset\Omega.
\]
Then $u\in C^\alpha(\Omega)$, and for every $\Omega'\Subset\Omega$,
\[
\sup_{\Omega'}|u|+
\sup_{\substack{x,y\in\Omega'\\x\ne y}}
\frac{|u(x)-u(y)|}{|x-y|^\alpha}
\leq c\bigl(M+\|u\|_{L^2(\Omega)}\bigr),
\]
where $c$ depends only on $n,\alpha,\Omega$, and $\Omega'$.
\end{corollary}
\begin{proof}
\sketch Poincar\'e's inequality gives
\[
\int_{B_r(x)}|u-u_{x,r}|^2
\leq c(n)r^2\int_{B_r(x)}|Du|^2
\leq c(n)M^2r^{n+2\alpha},
\]
so Theorem~\ref{hanlin:weaksol-part1-campanato} applies.
\end{proof}

\begin{lemma}\label{hanlin:weaksol-part1-lemma-3-3}\dcref{Lemma 3.3}\group{ch03-weak-solutions-part-i}
\source{han-lin-lecture-notes:page-0060}
Suppose $u\in H^1(\Omega)$ and, for some $\mu\in[0,n]$,
\[
\int_{B_r(x_0)}|Du|^2\leq Mr^\mu
\qquad\text{whenever }B_r(x_0)\subset\Omega.
\]
For every $\Omega'\Subset\Omega$ and every $B_r(x_0)\subset\Omega$ with
$x_0\in\Omega'$,
\[
\int_{B_r(x_0)}|u|^2
\leq c\left(M+\int_\Omega u^2\right)r^\lambda,
\]
where $\lambda=\mu+2$ if $\mu<n-2$, while $\lambda$ may be any number
in $[0,n]$ if $n-2\leq\mu<n$. The constant $c$ depends only on
$n,\lambda,\mu,\Omega$, and $\Omega'$.
\end{lemma}
\begin{proof}
\sketch For $R_0=\dist(\Omega',\partial\Omega)$, Poincar\'e's inequality
and the hypothesis give
\[
\int_{B_r(x_0)}|u-u_{x_0,r}|^2\leq c(n)Mr^{\mu+2}\leq c(n)Mr^\lambda.
\]
Consequently, for $0<\rho<r\leq R_0$,
\[
\int_{B_\rho(x_0)}u^2
\leq c(n)\left(\frac{\rho}{r}\right)^n
\int_{B_r(x_0)}u^2+Mr^\lambda.
\]
Lemma~\ref{hanlin:weaksol-part1-lemma-3-4} improves this to
\[
\int_{B_\rho(x_0)}u^2
\leq c\left(\frac{\rho}{r}\right)^\lambda
\int_{B_r(x_0)}u^2+M\rho^\lambda.
\]
Set $r=R_0$ and bound the remaining integral by $\int_\Omega u^2$.
\end{proof}

\begin{lemma}\label{hanlin:weaksol-part1-lemma-3-4}\dcref{Lemma 3.4}\group{ch03-weak-solutions-part-i}
\source{han-lin-lecture-notes:page-0061}
Let $\phi$ be nonnegative and nondecreasing on $[0,R]$. Suppose that
$A,B,\alpha,\beta\geq0$, $\beta<\alpha$, and
\[
\phi(\rho)\leq A\left(\left(\frac{\rho}{r}\right)^\alpha
+\epsilon\right)\phi(r)+Br^\beta
\qquad(0<\rho\leq r\leq R).
\]
For every $\gamma\in(\beta,\alpha)$ there is $\epsilon_0>0$ such that,
if $\epsilon<\epsilon_0$, then
\[
\phi(\rho)\leq c\left(
\left(\frac{\rho}{r}\right)^\gamma\phi(r)+Br^\beta\right)
\qquad(0<\rho\leq r\leq R),
\]
where $\epsilon_0$ and $c$ depend only on $A,\alpha,\beta$, and $\gamma$.
In particular,
\[
\phi(r)\leq c\left(\frac{\phi(R)}{R^\gamma}r^\gamma+Br^\beta\right)
\qquad(0<r\leq R).
\]
\end{lemma}
\begin{proof}
\sketch Choose $\tau\in(0,1)$ so that
$2A\tau^\alpha=\tau^\gamma$, and then choose
$\epsilon_0\tau^{-\alpha}<1$. Thus
\[
\phi(\tau r)\leq\tau^\gamma\phi(r)+Br^\beta.
\]
Iteration yields
\[
\phi(\tau^{k+1}r)
\leq\tau^{(k+1)\gamma}\phi(r)
+\frac{B\tau^{k\beta}r^\beta}{1-\tau^{\gamma-\beta}}.
\]
Selecting $k$ with $\tau^{k+2}r<\rho\leq\tau^{k+1}r$ proves both
estimates.
\end{proof}

\begin{theorem}\label{hanlin:weaksol-part1-john-nirenberg}\dcref{Theorem 3.5}\group{ch03-weak-solutions-part-i}
\source{han-lin-lecture-notes:page-0061,han-lin-lecture-notes:page-0062,han-lin-lecture-notes:page-0063}
Suppose $u\in L^1(\Omega)$ satisfies
\[
\int_{B_r(x)}|u-u_{x,r}|\leq Mr^n
\qquad\text{whenever }B_r(x)\subset\Omega.
\]
Then, for every $B_r(x)\subset\Omega$,
\[
\int_{B_r(x)}e^{\frac{p_0}{M}|u-u_{x,r}|}\leq Cr^n,
\]
where $p_0,C>0$ depend only on $n$.
\end{theorem}

\begin{proof}
\sketch Replace balls by comparable cubes and iterate
Lemma~\ref{hanlin:weaksol-part1-cz} on the oscillation. The resulting
geometric superlevel-set decay gives the exponential estimate.
\end{proof}

\begin{remark}\label{hanlin:weaksol-part1-bmo}\dcref{Remark 3.6}\group{ch03-weak-solutions-part-i}
\source{han-lin-lecture-notes:page-0061}
Functions satisfying the hypothesis of
Theorem~\ref{hanlin:weaksol-part1-john-nirenberg} have bounded mean
oscillation. The inclusion $L^\infty\subset\BMO$ is proper: on $(0,1)$,
$u(x)=\log x$ belongs to $\BMO$ but not to $L^\infty$.
\end{remark}

A dyadic subcube of the unit cube $Q_0$ is obtained by repeated subdivision
into $2^n$ congruent cubes; its predecessor is the cube from the preceding
generation that contains it.

\begin{lemma}\label{hanlin:weaksol-part1-cz}\dcref{Lemma 3.7}\group{ch03-weak-solutions-part-i}
\source{han-lin-lecture-notes:page-0062}
Let $f\in L^1(Q_0)$ be nonnegative and let
$\alpha>|Q_0|^{-1}\int_{Q_0}f$. There is a pairwise disjoint family of
dyadic cubes $\{Q_j\}$ in $Q_0$ such that
\[
f(x)\leq\alpha\quad\text{for a.e. }x\in Q_0\setminus\bigcup_jQ_j,
\]
and
\[
\alpha\leq\frac1{|Q_j|}\int_{Q_j}f<2^n\alpha.
\]
\end{lemma}
\begin{proof}
\sketch Retain every maximal dyadic cube whose average is at least $\alpha$
and subdivide every other cube. Maximality and the predecessor bound give the
two-sided average estimate. Outside the retained cubes, a nested sequence of
dyadic cubes has averages below $\alpha$ and diameters tending to zero, so the
Lebesgue differentiation theorem gives $f\leq\alpha$ almost everywhere.
\end{proof}

\section{H\"older Continuity of Solutions}

Assume that $a_{ij}\in L^\infty(B_1)$ is uniformly elliptic:
\[
\lambda|\xi|^2\leq a_{ij}(x)\xi_i\xi_j\leq\Lambda|\xi|^2
\qquad(x\in B_1,\ \xi\in\R^n),
\]
and that $u\in H^1(B_1)$ satisfies
\[
(\ast)\qquad
\int_{B_1}a_{ij}D_i uD_j\varphi+cu\varphi
=\int_{B_1}f\varphi
\qquad(\varphi\in H^1_0(B_1)).
\]

\begin{theorem}\label{hanlin:weaksol-part1-holder-sol}\dcref{Theorem 3.8}\group{ch03-weak-solutions-part-i}
\source{han-lin-lecture-notes:page-0064,han-lin-lecture-notes:page-0065,han-lin-lecture-notes:page-0066,han-lin-lecture-notes:page-0067,han-lin-lecture-notes:page-0068,han-lin-lecture-notes:page-0069,han-lin-lecture-notes:page-0070}
Let $u\in H^1(B_1)$ solve $(\ast)$. If
$a_{ij}\in C^0(\overline{B_1})$, $c\in L^n(B_1)$, and
$f\in L^q(B_1)$ for some $q\in(n/2,n)$, then
$u\in C^\alpha(B_1)$ for $\alpha=2-n/q\in(0,1)$.
Moreover, there is $R_0>0$ such that, for $x\in B_{1/2}$ and $r\leq R_0$,
\[
\int_{B_r(x)}|Du|^2
\leq Cr^{n-2+2\alpha}
\bigl(\|f\|_{L^q(B_1)}^2+\|u\|_{H^1(B_1)}^2\bigr).
\]
Here $R_0$ and $C$ depend only on $\lambda,\Lambda,\tau$, and
$\|c\|_{L^n}$, where
\[
|a_{ij}(x)-a_{ij}(y)|\leq\tau(|x-y|)
\qquad(x,y\in B_1).
\]
If $c\equiv0$, $\|u\|_{H^1(B_1)}$ may be replaced by
$\|Du\|_{L^2(B_1)}$.
\end{theorem}

\begin{proof}
\sketch Decompose into a frozen-coefficient harmonic part and an error.
Energy estimates, Corollary~\ref{hanlin:weaksol-part1-cor-3-10}, and the
Campanato iteration yield the Morrey estimate and H\"older continuity.
\end{proof}

\begin{lemma}\label{hanlin:weaksol-part1-harmonic-lemma}\dcref{Lemma 3.9}\group{ch03-weak-solutions-part-i}
\source{han-lin-lecture-notes:page-0064}
Let $(a_{ij})$ be constant and satisfy
\[
\lambda|\xi|^2\leq a_{ij}\xi_i\xi_j\leq\Lambda|\xi|^2
\qquad(\xi\in\R^n).
\]
If $w\in H^1(B_r(x_0))$ is a weak solution of
\[
a_{ij}D_iD_jw=0\qquad\text{in }B_r(x_0),
\]
then, for $0<\rho\leq r$,
\[
\int_{B_\rho(x_0)}|Dw|^2
\leq C\left(\frac\rho r\right)^n\int_{B_r(x_0)}|Dw|^2,
\]
and
\[
\int_{B_\rho(x_0)}|Dw-(Dw)_{x_0,\rho}|^2
\leq C\left(\frac\rho r\right)^{n+2}
\int_{B_r(x_0)}|Dw-(Dw)_{x_0,r}|^2,
\]
where $C$ depends only on $n$ and $\Lambda/\lambda$.
\end{lemma}
\begin{proof}
\sketch Every derivative of $w$ solves the same constant-coefficient equation;
the two estimates follow by applying the harmonic derivative estimate of
Lemma~\ref{hanlin:lemma-constant-coeff-compactness} to $Dw$.
\end{proof}

\begin{corollary}\label{hanlin:weaksol-part1-cor-3-10}\dcref{Corollary 3.10}\group{ch03-weak-solutions-part-i}
\source{han-lin-lecture-notes:page-0065}
Suppose $w$ is as in Lemma~\ref{hanlin:weaksol-part1-harmonic-lemma} and
$u\in H^1(B_r(x_0))$. Then, for
$0<\rho\leq r$,
\[
\int_{B_\rho(x_0)}|Du|^2
\leq C\left\{
\left(\frac\rho r\right)^n\int_{B_r(x_0)}|Du|^2
+\int_{B_r(x_0)}|D(u-w)|^2\right\},
\]
and
\[
\int_{B_\rho(x_0)}|Du-(Du)_{x_0,\rho}|^2
\leq C\left\{
\left(\frac\rho r\right)^{n+2}
\int_{B_r(x_0)}|Du-(Du)_{x_0,r}|^2
+\int_{B_r(x_0)}|D(u-w)|^2\right\},
\]
where $C$ depends only on $n$ and $\Lambda/\lambda$.
\end{corollary}
\begin{proof}
\sketch Write $v=u-w$. Apply
Lemma~\ref{hanlin:weaksol-part1-harmonic-lemma} to $w=u-v$ and use
$|a+b|^2\leq2|a|^2+2|b|^2$, first for $Du$, then after subtracting the
corresponding averages.
\end{proof}

\section{H\"older Continuity of Gradients}

Continue to assume uniform ellipticity and $(\ast)$ as above.

\begin{theorem}\label{hanlin:weaksol-part1-grad-holder}\dcref{Theorem 3.11}\group{ch03-weak-solutions-part-i}
\source{han-lin-lecture-notes:page-0068,han-lin-lecture-notes:page-0069,han-lin-lecture-notes:page-0070,han-lin-lecture-notes:page-0071}
Let $u\in H^1(B_1)$ solve $(\ast)$. Suppose
$a_{ij}\in C^\alpha(B_1)$ and $c,f\in L^q(B_1)$ for some $q>n$, where
$\alpha=1-n/q\in(0,1)$. Then $Du\in C^\alpha(B_1)$. Moreover, there is
$R_0\in(0,1)$ such that, for $x\in B_{1/2}$ and $r\leq R_0$,
\[
\int_{B_r(x)}|Du-(Du)_{x,r}|^2
\leq Cr^{n+2\alpha}
\left(\|f\|_{L^q(B_1)}^2+\|u\|_{H^1(B_1)}^2\right),
\]
where $R_0$ and $C$ depend only on $\lambda$,
$|a_{ij}|_{C^\alpha}$, and $|c|_{L^q}$.
\end{theorem}
\begin{proof}
\sketch Use the frozen-coefficient decomposition $u=v+w$ from the proof of
Theorem~\ref{hanlin:weaksol-part1-holder-sol}.
Corollary~\ref{hanlin:weaksol-part1-cor-3-10} gives, for
$0<\rho\leq r$,
\begin{equation}
\begin{aligned}
\int_{B_\rho(x_0)}|Du-(Du)_{x_0,\rho}|^2
\leq C\biggl\{&
\left(\frac\rho r\right)^{n+2}
\int_{B_r(x_0)}|Du-(Du)_{x_0,r}|^2
+\tau^2(r)\int_{B_r(x_0)}|Du|^2\\
&+\left(\int_{B_r(x_0)}|c|^n\right)^{2/n}
\int_{B_r(x_0)}u^2
+\left(\int_{B_r(x_0)}|f|^{2n/(n+2)}\right)^{(n+2)/n}
\biggr\}.
\end{aligned}
\end{equation}
Since $\alpha=1-n/q$,
\[
\left(\int_{B_r}|f|^{2n/(n+2)}\right)^{(n+2)/n}
\leq r^{n+2\alpha}\|f\|_{L^q(B_1)}^2,
\qquad
\left(\int_{B_r}|c|^n\right)^{2/n}
\leq r^{2\alpha}\|c\|_{L^q(B_1)}^2.
\]
For constant $a_{ij}$ and $c=0$,
Lemma~\ref{hanlin:weaksol-part1-lemma-3-4} applied to the preceding estimate
gives
the required Campanato estimate. If $c=0$ and
$\tau(r)\leq Cr^\alpha$, the companion energy estimate first yields, for
every small $\delta>0$,
\[
\int_{B_r(x_0)}|Du|^2
\leq Cr^{n-2\delta}
\left(\|f\|_{L^q(B_1)}^2+\|Du\|_{L^2(B_1)}^2\right).
\]
Substitution in the oscillation estimate gives
$Du\in C^{\alpha-\delta}_{\mathrm{loc}}$, hence a local $L^\infty$ bound
for $Du$. Substituting that bound back into the same estimate and
applying Lemma~\ref{hanlin:weaksol-part1-lemma-3-4} gives exponent $\alpha$.

In the general case, the same bootstrap, with
Lemma~\ref{hanlin:weaksol-part1-lemma-3-3} controlling
$\int_{B_r}u^2$, yields
\[
\sup_{B_{3/4}}|u|^2+\sup_{B_{3/4}}|Du|^2
\leq C\left(\|f\|_{L^q(B_1)}^2+\|u\|_{H^1(B_1)}^2\right).
\]
The oscillation estimate then reduces to
\[
\int_{B_\rho(x_0)}|Du-(Du)_{x_0,\rho}|^2
\leq C\left\{
\left(\frac\rho r\right)^{n+2}
\int_{B_r(x_0)}|Du-(Du)_{x_0,r}|^2
+r^{n+2\alpha}
\left(\|f\|_{L^q}^2+\|u\|_{H^1}^2\right)\right\}.
\]
Lemma~\ref{hanlin:weaksol-part1-lemma-3-4} and
Theorem~\ref{hanlin:weaksol-part1-campanato} complete the proof.
\end{proof}

\begin{example}\label{hanlin:weaksol-part1-example-3-12}\dcref{Example 3.12}\group{ch03-weak-solutions-part-i}
\source{han-lin-lecture-notes:page-0072}
Let $B_R\subset\R^n$, $R<1$, and set
\[
u(x)=(x_1^2-x_2^2)(-\log|x|)^{1/2}.
\]
Then
\[
\Delta u=
\frac{x_2^2-x_1^2}{2|x|^2}
\left(\frac{n+2}{(-\log|x|)^{1/2}}
+\frac{1}{2(-\log|x|)^{3/2}}\right),
\]
whose right-hand side extends continuously by zero at the origin, and
$u=\sqrt{-\log R}(x_1^2-x_2^2)$ on $\partial B_R$. The function $u$ is a
weak solution and is smooth away from the origin, but
$\lim_{|x|\to0}D_{11}u=\infty$, so $u\notin C^2(B_R)$. No classical
solution with these boundary data exists: if $v$ were one, then $u-v$ would
be bounded and harmonic on $B_R\setminus\{0\}$, hence would extend
harmonically across the origin by
Theorem~\ref{hanlin:thm-removable-singularity}, contradicting the divergence
of $D_{11}u$.
\end{example}
